\documentclass[10pt,a4paper]{amsart}
\usepackage[margin=19mm]{geometry}
\usepackage{amsmath,amssymb,mathtools}
\usepackage[scr=rsfs,cal=euler]{mathalfa}
\usepackage{xfrac}
\usepackage[unicode,hidelinks,bookmarks=false]{hyperref}
\newcommand{\ie}{i.e.\ }
\newcommand{\cf}{cf.\ }
\newcommand{\resp}{respectively\ }
\newcommand{\Subsection}{\subsection}
\providecommand{\nobreakdash}{\nobreak\hbox{-}}
\setlength{\emergencystretch}{2em}
\multlinegap0pt
\usepackage{bm}
\usepackage{bbm}
\usepackage{dsfont}
\usepackage{xspace}
\usepackage{cancel}
\usepackage{mathtools}
\usepackage{cleveref}
\usepackage{amsthm}
\makeatletter
\@ifundefined{theorem}{\newtheorem{theorem}{Theorem}}{}
\@ifundefined{definition}{\newtheorem{definition}[theorem]{Definition}}{}
\@ifundefined{proposition}{\newtheorem{proposition}[theorem]{Proposition}}{}
\makeatother
\title[Limits of Convergence-Rate Control for Open-Weight Safety]{Limits of Convergence-Rate Control for Open-Weight Safety}
\author{Domenic Rosati, Xijie Zeng, Hong Huang, Sebastian Dionicio, Subhabrata Majumdar, Frank Rudzicz,, Hassan Sajjad}
\date{Source: arXiv:2602.18868v1 (CC BY 4.0)}
\begin{document}
\maketitle

\noindent\emph{Source and licence.} Limits of Convergence-Rate Control for Open-Weight Safety. Version arXiv:2602.18868v1 (CC BY 4.0), distributed under the Creative Commons Attribution 4.0 International licence, as recorded in the arXiv licence metadata for this exact version and shown on the abstract page https://arxiv.org/abs/2602.18868v1. Full rights notice: licenses/arxiv-2602.18868-CC-BY-4.0.txt.\par\smallskip

\noindent\textit{Editorial scope.} Counted unit: the Theorem whose source label is \texttt{thm:tensor-b3} in the article (source0.tex lines 2983--2992, source label \texttt{thm:tensor-b3}), delivered together with the article's own complete original proof of it (source0.tex lines 2993--3068, one \texttt{proof} environment of 76 lines). No mathematics is written, repaired or re-derived: statement and proof are the article's own text line for line. Required context retained uncounted: prop:L-lower-bound-app (lines 787--791); def:lip-cont-hes (lines 2620--2629); def:dft (lines 2879--2900); thm:tsvd (lines 2937--2975). Every definition, display and label of those lines is retained unchanged. Omitted from this package and not redistributed: 1--786, 792--2619, 2630--2878, 2901--2936, 2976--2982, 3069--3716. Nothing inside a retained excerpt is omitted or altered: every retained body is delivered exactly as the article's lines are written, so no omission receipt of any kind accompanies this package. Citations: the retained excerpts carry no citation command at all, so no bibliography entry of the article is transcribed and no reference is redistributed. Reference adaptation: none was needed and none was made; every reference inside the delivered unit resolves inside it, so no unresolved reference or citation is produced. Printed slips kept literal: no printed word, symbol, hypothesis or number of the counted statement or of its proof was altered. Layout and class: the article's own class is \texttt{article} with packages including microtype, graphicx, subcaption, booktabs, hyperref, icml2026, amsmath, amssymb, mathtools, amsthm, cleveref, backref, stfloats, siunitx; the wrapper is the lane's pinned \texttt{amsart} editorial wrapper at 10pt on A4, which restates the theorem-like environments the retained text needs and adds the article's own macro definitions that the retained text needs (none), with extra wrapper packages (bm, bbm, dsfont, xspace, cancel, mathtools, cleveref). The delivered numbering is the unit's own. Assets: no retained span contains a figure environment, a graphics-inclusion command, a PDF-page inclusion, a table or a TikZ picture; no image file is copied into this package, the package declares no asset dependency and no image file is redistributed with it. Licence: version arXiv:2602.18868v1 of this article is distributed under the Creative Commons Attribution 4.0 International licence, as recorded in the arXiv licence metadata for this exact version and shown on the abstract page https://arxiv.org/abs/2602.18868v1. A full rights notice is delivered at licenses/arxiv-2602.18868-CC-BY-4.0.txt. Characters: every retained excerpt is pure ASCII, so the delivered engine, XeLaTeX with its default Latin Modern text font, reports no missing character.
\begin{proposition}
\label{prop:L-lower-bound-app}
The smoothness quantifier $L \in \mathbb{R}$ for the function $\mathcal{L}$ is greater than or equal to the maximum singular value of the Hessian $H$ of that function. In other words the following formula holds
$L \geq \sigma_1(H_{\theta}^{\mathcal{L}})$.
\end{proposition}

\begin{definition}[$L_2$-Lipschitz Continuous Hessian]
\label{def:lip-cont-hes}
A twice continuously differentiable function $f:\mathbb{R}^n\to\mathbb{R}$ 
is said to have an \emph{$L_2$-Lipschitz continuous Hessian} if there exists a constant $L_2>0$ such that
for all $x,y\in\mathbb{R}^n$,
\[
\|\nabla^{2} f(x)-\nabla^{2} f(y)\|\le L_2\|x-y\|.
\]
in which case we call $L_2$ the \emph{Hessian Lipschitz constant}.
\end{definition}

\begin{definition}[Discrete Fourier Transform along the third dimension]
\label{def:dft}
For a tensor $\mathcal{A}\in\mathbb{R}^{m\times n\times p}$, 
the Discrete Fourier Transform (DFT) along the third dimension is defined as
\[
\widehat{\mathcal{A}}(:,:,k)
= \sum_{t=1}^{p}\mathcal{A}(:,:,t)\,e^{-2\pi i (t-1)(k-1)/p}, 
\quad k=1,\dots,p.
\]
The inverse DFT is given by
\[
\mathcal{A}(:,:,t)
= \frac{1}{p}\sum_{k=1}^{p}\widehat{\mathcal{A}}(:,:,k)\,e^{2\pi i (t-1)(k-1)/p}.
\]

Under this transform, the block-circulant matrix of $\mathcal{A}$ is block-diagonalized as
\[
(F_p\otimes I_m)\,\mathrm{bcirc}(\mathcal{A})\,(F_p^{-1}\otimes I_n)
= \mathrm{diag}\big(\widehat{\mathcal{A}}(:,:,1),\dots,\widehat{\mathcal{A}}(:,:,p)\big),
\]
where $F_p$ is the $p\times p$ DFT matrix with entries $[F_p]_{jk} = e^{-2\pi i (j-1)(k-1)/p}$, and $\otimes$ denotes the Kronecker product.
\end{definition}

\begin{theorem}[Tensor Singular Value Decomposition (T-SVD) and Tensor Singular Values]
\label{thm:tsvd}
Let $\mathcal{A}\in\mathbb{R}^{m\times n\times p}$ be a third-order tensor. 
Then $\mathcal{A}$ admits the following factorization under the T-product:
\[
\mathcal{A} = \mathcal{U} * \mathcal{S} * \mathcal{V}^\top,
\]
where $\mathcal{U}\in\mathbb{R}^{m\times m\times p}$ and 
$\mathcal{V}\in\mathbb{R}^{n\times n\times p}$ are orthogonal tensors satisfying 
$\mathcal{U}^\top * \mathcal{U} = \mathcal{I}$ and $\mathcal{V}^\top * \mathcal{V} = \mathcal{I}$, 
and $\mathcal{S}\in\mathbb{R}^{m\times n\times p}$ is an \emph{f-diagonal tensor}, 
that is, each frontal slice $\mathcal{S}(:,:,k)$ is diagonal.
According to Definition~\ref{def:dft}, if 
$\widehat{\mathcal{A}}=\mathrm{fft}(\mathcal{A},[],3)$, 
then each frontal slice of $\widehat{\mathcal{A}}$ admits a standard matrix SVD:
\[
\widehat{\mathcal{A}}(:,:,k)
= \widehat{U}_k\,\widehat{S}_k\,\widehat{V}_k^\top, 
\quad k=1,\dots,p,
\]
where $\widehat{U}_k,\widehat{V}_k$ are orthogonal matrices and 
$\widehat{S}_k=\mathrm{diag}\big(\sigma_1^{(k)},\dots,\sigma_{r}^{(k)}\big)$ with $r=\min(m,n)$. Applying the inverse transform defined in Definition~\ref{def:dft} yields
\[
\mathcal{U}=\mathrm{ifft}(\widehat{\mathcal{U}},[],3), \quad
\mathcal{S}=\mathrm{ifft}(\widehat{\mathcal{S}},[],3), \quad
\mathcal{V}=\mathrm{ifft}(\widehat{\mathcal{V}},[],3).
\]
The \bf{tensor singular values} of $\mathcal{A}$ are defined as the diagonal entries of the first frontal slice $\mathcal{S}(:,:,1)$, denoted by
\[
\varsigma_i(\mathcal{A}) := \mathcal{S}(i,i,1)
= \frac{1}{p}\sum_{k=1}^{p}\sigma_i(\widehat{\mathcal{A}}(:,:,k)), 
\quad i=1,\dots,r,
\]
Hence, the tensor singular values satisfy the monotonicity
\[
\varsigma_1(\mathcal{A}) \ge \varsigma_2(\mathcal{A}) \ge \cdots 
\ge \varsigma_r(\mathcal{A}) \ge 0.
\]
\end{theorem}

\begin{theorem}[Tensor Version of \cref{prop:L-lower-bound-app}]
\label{thm:tensor-b3}
Let $f:\mathbb{R}^n\to\mathbb{R}$ be thrice continuously differentiable with an
$L_2$-Lipschitz continuous Hessian as defined in
Definition~\ref{def:lip-cont-hes}. Let $\mathcal{T}(x)=\nabla^3 f(x)\in\mathbb{R}^{n\times n\times n}$ denote the third-order derivative tensor.
Then the largest tensor singular value of $\mathcal{T}(x)$ under the T--SVD framework satisfies
\[
\varsigma_1\!\big(\mathcal{T}(x)\big) \;\le\; L_2.
\]
\end{theorem}

\begin{proof}
Let $f$ satisfy the $L_2$-Lipschitz Hessian condition in 
Definition~\ref{def:lip-cont-hes}, namely
\[
\|\nabla^{2} f(u)-\nabla^{2} f(v)\|\le L_2\|u-v\|,
\qquad \forall\,u,v\in\mathbb{R}^n.
\]
Fix a point $x\in\mathbb{R}^n$ and a unit direction $h\in\mathbb{R}^n$ with 
$\|h\|_2=1$.  
Define the matrix-valued function
\[
G(t) := \nabla^2 f(x + t h), \qquad t\in [0,1].
\]
Then by Definition~\ref{def:lip-cont-hes},
\[
\|G(t)-G(s)\|
= \|\nabla^2 f(x+th)-\nabla^2 f(x+sh)\|
\le L_2|t-s|.
\]
Hence $G$ is an $L_2$-Lipschitz function of $t$.  
Since $f$ is $C^3$ and $G$ is differentiable then
\[
G'(t)
= \frac{d}{dt}\nabla^2 f(x+th)
= \nabla^3 f(x+th)[\,\cdot,\cdot,h\,].
\]
Since $G$ is differentiable and its derivative $G'(t)$ is continuous,
for each fixed $t$ we have
\[
\|G'(t)\|
= \Bigl\|\lim_{s\to t}\frac{G(t)-G(s)}{t-s}\Bigr\|
\le \limsup_{s\to t}\frac{\|G(t)-G(s)\|}{|t-s|}
\le L_2, \qquad  \forall\,t\in[0,1].
\]
Therefore,
\begin{equation}\label{eq:g.9.1}
\big\|\nabla^3 f(x+th)[\,\cdot,\,\cdot,\,h\,]\big\|_2
= \|G'(t)\|
\le L_2,
\qquad \forall\,t\in[0,1],\ \|h\|_2=1.
\end{equation}
Applying the discrete Fourier transform along the third tensor dimension 
(Definition~\ref{def:dft}) gives
\[
\widehat{\mathcal{T}}
= \mathrm{fft}(\mathcal{T}(x),[],3),
\qquad 
\mathcal{T}(x)=\nabla^3 f(x).
\]
Its $k$-th frontal slice satisfies
\[
\widehat{\mathcal{T}}(:,:,k)
=\mathcal{T}(x)[\,\cdot,\cdot,\phi_k\,],
\]
where $\phi_k$ denotes the $k$-th unit-norm Fourier basis vector under the unitary DFT convention adopted throughout this appendix. 
% fix comment 1
Moreover, the bound in \eqref{eq:g.9.1} extends to complex unit vectors by identifying $\mathbb{C}^n$ with $\mathbb{R}^{2n}$ and using the same operator norm.

Since $\|\phi_k\|_2=1$, substituting $h=\phi_k$ into 
\eqref{eq:g.9.1} yields
\[
\sigma_{1}\!\big(\widehat{\mathcal{T}}(:,:,k)\big) =
\big\|\widehat{\mathcal{T}}(:,:,k)\big\|_2
\le L_2,
\qquad k=1,\dots,p.
\]
From Theorem~\ref{thm:tsvd}, the largest tensor singular values of $\mathcal{T}(x)$ are
\[
\varsigma_1\!\big(\mathcal{T}(x)\big)
= \frac{1}{p}\sum_{k=1}^{p}\sigma_{1}\!\big(\widehat{\mathcal{T}}(:,:,k)\big)
\le \frac{1}{p}\sum_{k=1}^{p}L_2
= L_2.
\]
Therefore the largest tensor singular value of the third-order derivative tensor
$\nabla^3 f(x)$ is bounded above by the Hessian Lipschitz constant $L_2$.
\end{proof}

\end{document}
